
%\section{Naive Approach}

%\begin{frame}{Diagonal Equation Annotation}
%	
%	% Define an explanation node style
%	\tikzstyle{annot} = [font=\small, text=red, align=left]
%	
%	\begin{equation*}
%		% Mark a specific part of your equation using \tikzmarknode
%		\tikzmarknode{termA}{A} + B = C
%	\end{equation*}
%	
%	\begin{tikzpicture}[overlay, remember picture]
%		% Draw a diagonal arrow/line from the term to an offset coordinate and place text
%		\draw[->, thick, red] (termA.south west) -- ++(-0.5cm, -0.1cm) 
%		node[annot, below left] {Diagonal explanation text here};
%	\end{tikzpicture}
%	
%\end{frame}

\begin{frame}{Coarse Approach 1}
	\vspace{-0.5cm}
	\begin{align*}
		\begin{bNiceMatrix}
			x[0]\\
			x[1]\\
			x[2]\\
			\vdots
		\end{bNiceMatrix} &=
		\begin{bmatrix}
			1 & 1 & 1 & \cdots & 1\\
			1 & e^{i\frac{2\pi}{N}} & e^{i\frac{4\pi}{N}} & \cdots & e^{i\frac{2\pi(N-1)}{N}}\\
			1 & e^{i\frac{4\pi}{N}} & e^{i\frac{8\pi}{N}} & \cdots & e^{i\frac{4\pi(N-1)}{N}}\\
			\vdots & \vdots & \vdots & \ddots & \vdots\\
			1 & e^{i\frac{2\pi(N-1)}{N}} &
			e^{i\frac{4\pi(N-1)}{N}} &
			\cdots &
			e^{i\frac{2\pi(N-1)^2}{N}}
		\end{bmatrix}
		\begin{bNiceMatrix}
		X[0]\\
		X[1]\\
		X[2]\\
		\vdots\\
		X[N-1]
		\end{bNiceMatrix}
	\end{align*}
\end{frame}

\begin{frame}[fragile]{Coarse Approach 1}
	Compute correlation surface from $Nsamples=256$ complex coefficients. Live-compute twiddles. $O(n^2)$ complex multiplications
	
	\begin{lstlisting}[basicstyle=\ttfamily\scriptsize]
auto [coeffs_a_conj, coeffs_b] = // 256 coefficients each
array_r<Nsamples> surface;

auto time_start = time_func();
// loop is N*N = 256*256 = 65536 multiplications
for (size_t out_idx = 0; out_idx < Nsamples; ++out_idx)
{
	Real sum(0.0);
	for (size_t i = 0; i < coeffs_a_conj.size(); ++i)
	{
		auto a_conj = coeffs_a_conj[i];
		auto b = coeffs_b[i];
		
		Real freq = output_index * i / Real(domain::WindowSize);
		Complex term_i = a_conj * b * std::exp(twopi * freq);
		sum += term_i.real();
	}
	surface[out_idx] = sum;
}
auto time_stop = time_func();
	\end{lstlisting}
	
\end{frame}

\begin{frame}{Results}
	\coarsetable{1}{1}
\end{frame}


\begin{frame}{Coarse Approach 2}
	\vspace{-0.5cm}
	\begin{align*}
		\begin{bNiceMatrix}
			x[0]\\
			x[1]\\
			x[2]\\
			\vdots
		\end{bNiceMatrix} &=
		\begin{bNiceMatrix}
			\Block[fill=red!20]{9-1}{} 1 & 1 & \cdots & \Block[fill=red!20]{9-1}{} 1 & \Block[fill=yellow!20]{9-2}{} \cdots & 1\\
			1 & e^{i\frac{2\pi}{N}} & \cdots & -1 & \cdots & e^{i\frac{2\pi(N-1)}{N}}\\
			1 & e^{i\frac{4\pi}{N}} & \cdots & 1 & \cdots & e^{i\frac{4\pi(N-1)}{N}}\\
			\vdots & \vdots & \vdots & \vdots & \ddots & \vdots\\
			1 & e^{i\frac{2\pi(\frac{N}{2}-1)}{N}} & \cdots & -1 & \cdots & e^{i\frac{2\pi(\frac{N}{2}-1)(N-1)}{N}}\\
			1 & e^{i\frac{2\pi(\frac{N}{2})}{N}} & \cdots & 1 & \cdots &  e^{i\frac{2\pi(\frac{N}{2})(N-1)}{N}}\\
			1 & e^{i\frac{2\pi(\frac{N}{2}+1)}{N}} & \cdots & -1 & \cdots & e^{i\frac{2\pi(\frac{N}{2}+1)(N-1)}{N}}\\
			\vdots & \vdots & \vdots & \vdots & \ddots & \vdots\\
			1  &  e^{i\frac{2\pi(N-1)}{N}} &\cdots & -1 & \cdots & e^{i\frac{2\pi(N-1)^2}{N}}
		\end{bNiceMatrix}
		\begin{bNiceMatrix}
			\Block[fill=red!20]{1-1}{} \tikzmarknode{termDC}{X[0]}\\
			X[1]\\
			X[2]\\
			\vdots\\
			X[\frac{N}{2}-1] \\
			\Block[fill=red!20]{1-1}{} \tikzmarknode{termNY}{X[\frac{N}{2}]} \\
			\Block[fill=yellow!20]{3-1}{} X[\frac{N}{2}+1]\\
			\vdots \\
			\tikzmarknode{termNeg}{X[N-1]}
		\end{bNiceMatrix}
	\end{align*}
	
	
	\begin{tikzpicture}[overlay, remember picture]
		% Draw a diagonal arrow/line from the term to an offset coordinate and place text
		\draw[<-, thick, red] (termDC.west) -- ++(-0.5cm, 1.26cm) 
		node[annot, left] {DC Coefficient (freq=0 or +N or -N)};
	\end{tikzpicture}
	
	\begin{tikzpicture}[overlay, remember picture]
		% Draw a diagonal arrow/line from the term to an offset coordinate and place text
		\draw[<-, thick, red] (termNY.west) -- ++(-1.0cm, 3.5cm) 
		node[annot, left] {Nyquist Coefficient (freq=+N/2 or -N/2)};
	\end{tikzpicture}
	
	\begin{tikzpicture}[overlay, remember picture]
		% Draw a diagonal arrow/line from the term to an offset coordinate and place text
		\draw[<-, thick, red] (termNeg.south) -- ++(-0.1cm, -1.0cm) 
		node[annot, left] {Coefficients for negative freqs. complex conjugates of the positive freqs};
	\end{tikzpicture}
\end{frame}

\begin{frame}[fragile]{Coarse Approach 2}
	Exploit conjugate symmetry to compute correlation surface using $\frac{Nsamples}{2}+1=129$ complex coefficients . Live-compute twiddles. $O(n^2)$ complex multiplications. Multiply the terms with conjugate symmetry by 2 and take their real component since $z + \bar{z} = 2\operatorname{Re}(z) $
	
	
	\begin{lstlisting}[basicstyle=\ttfamily\scriptsize]
auto [coeffs_a_conj, coeffs_b] = // 129 coefficients each
array_r<Nsamples> surface;
		
for (size_t i=1; i < coeffs_a_conj.size()-1; ++i)
{
	coeffs_a_conj[i] *= 2;
}
		
auto time_start = time_func();
// same nested loop structure. 
// Now N*(N/2+1) = 256*129 = 33024 mults
auto time_stop = time_func();
	\end{lstlisting}
	
\end{frame}


\begin{frame}{Results}
	\coarsetable{1}{2}
\end{frame}


\begin{frame}{Coarse Approach 3}
	\vspace{-0.5cm}
	\begin{align*}
		\begin{bNiceMatrix}
			x[0]\\
			x[1]\\
			x[2]\\
			\vdots
		\end{bNiceMatrix} &=
		\begin{bNiceMatrix}
			1 & 2 & 2 & \Block[fill=red!20]{9-1}{} \cancel{\cdots} & 1 & \cancel{\cdots} \\
			1 & 2e^{i\frac{2\pi}{N}} & 2e^{i\frac{4\pi}{N}} & \cancel{\cdots} & -1 & \cancel{\cdots} \\
			1 & 2e^{i\frac{4\pi}{N}} & 2e^{i\frac{8\pi}{N}} & \cancel{\cdots} & 1 & \cancel{\cdots} \\
			\vdots & \vdots & \vdots & \vdots & \ddots & \vdots\\
			1 & 2e^{i\frac{2\pi(\frac{N}{2}-1)}{N}} & 2e^{i\frac{4\pi(\frac{N}{2}-1)}{N}} & \cancel{\cdots} & -1 & \cancel{\cdots} \\
			1 & 2e^{i\frac{2\pi(\frac{N}{2})}{N}} & 2e^{i\frac{4\pi(\frac{N}{2})}{N}} & \cancel{\cdots} & 1 & \cancel{\cdots} \\
			1 & 2e^{i\frac{2\pi(\frac{N}{2}+1)}{N}}  & 2e^{i\frac{4\pi(\frac{N}{2}+1)}{N}}& \cancel{\cdots} & -1 & \cancel{\cdots} \\
			\vdots & \vdots & \vdots & \vdots & \vdots & \ddots \\
			1  &  2e^{i\frac{2\pi(N-1)}{N}} & 2e^{i\frac{4\pi(N-1)}{N}} & \cancel{\cdots} & -1 & \cancel{\cdots} 
		\end{bNiceMatrix}
		\begin{bNiceMatrix}
			X[0]\\
			X[1]\\
			X[2]\\
			\Block[fill=red!20]{1-1}{} \tikzmarknode{termsNearZero}{\vdots}\\
			X[\frac{N}{2}] \\
			\tikzmarknode{termNeg}{\vdots}
		\end{bNiceMatrix}
	\end{align*}
	
	
	
	\begin{tikzpicture}[overlay, remember picture]
		% Draw a diagonal arrow/line from the term to an offset coordinate and place text
		\draw[<-, thick, red] (termsNearZero.north west) -- ++(-1.0cm, 3.0cm) 
		node[annot, left] {Most coefficients $\approx 0$ from chirp)};
	\end{tikzpicture}
\end{frame}


\begin{frame}[fragile]{Coarse Approach 3}
	Compute correlation surface using $N=60$ complex coefficients by omitting coefficients near $0$ and exploiting conjugate symmetry. Live compute twiddles. $O(nm)$
	
	\begin{lstlisting}[basicstyle=\ttfamily\scriptsize]
auto [coeffs_a_conj, coeffs_b] = // 129 coefficients each
array_r<Nsamples> surface;

for (size_t i=1; i < coeffs_a_conj.size()-1; ++i)
{
	coeffs_a_conj[i] *= 2;
}

// count the first M terms with magnitude larger than a threshold

auto time_start = time_func();
// same nested loop structure. Now N*M = 256*60 = 15360 mults
auto time_stop = time_func();
	\end{lstlisting}
	
\end{frame}

\begin{frame}{Approach 3 Not Viable}
	\begin{itemize}
		\item The terms that were near zero were not identically zero. 
		\item This introduced an error of $\approx$ 0.1 on a scale of 1E3 - 1E4)
		\item Previous errors were \textless 1E-3
	\end{itemize}
\end{frame}

\begin{frame}{Results}
	\coarsetable{1}{3}
\end{frame}


\begin{frame}[fragile]{Coarse Approach 4}
	Compute correlation surface using $Nsamples=256$ complex coefficients by exploiting conjugate symmetry. precompute twiddles. $O(n^2)$
	
	\begin{lstlisting}[basicstyle=\ttfamily\scriptsize, 
		linebackgroundcolor={%
			\ifnum\value{lstnumber}=10 \color{yellow!30}\fi
		}]
auto time_start = time_func();

constexpr size_t Noutputs_used = Nsamples/2;
for (size_t out_idx = 0; out_idx < Noutputs_used; ++out_idx)
{
	Real sum(0.0);
	for (size_t i = 0; i < Ncoeffs; ++i)
	{
		auto b = coeffs_b[i];
		Complex term_i =  b * Wn[out_idx][i]; // a_conj baked in
		sum += term_i.real();
	}
	surface[out_idx] = sum;
}
auto time_stop = time_func();
	\end{lstlisting}
	
\end{frame}


\begin{frame}{Results}
	\coarsetable{1}{4}
\end{frame}


\begin{frame}[fragile]{Coarse Approach 5}
	A Fourier Transform can be computed recursively. This is Radix-2:
	
	\begin{align}
		\mathbf{X} &= \mathcal{F}\{\mathbf{x}\}  \\
		X[k] &= \sum_{n=0}^{N-1} x[n] e^{-i\frac{2\pi}{N}kn}, \qquad k=0,1,\ldots,N-1 \\
		&= \sum_{n=0}^{N/2-1} x[n] e^{-i\frac{2\pi}{N}k(2n)} + \sum_{n=0}^{N/2-1} x[n] e^{-i\frac{2\pi}{N}k(2n+1)} \\
		&= \sum_{n=0}^{N/2-1} x[n] e^{-i\frac{2\pi}{N/2}kn} + e^{-i\frac{2\pi}{N}kn} \sum_{n=0}^{N/2-1} x[n] e^{-i\frac{2\pi}{N/2}kn} \\
		X[k] &= \mathcal{F}\{\mathbf{x}_{\texttt{evens}}\}  + e^{-i\frac{2\pi}{N}kn} \mathcal{F}\{\mathbf{x}_{\texttt{odds}}\} \\
		X[k+\frac{N}{2}] &= \mathcal{F}\{\mathbf{x}_{\texttt{evens}}\}  - e^{-i\frac{2\pi}{N}kn} \mathcal{F}\{\mathbf{x}_{\texttt{odds}}\}
	\end{align}
	turtles all the way down
	
\end{frame}

\begin{frame}
	\begin{itemize}
		\item You can radix by any prime number, but N must be a multiple/power of that number
		\item Radix-2, Radix-3, Radix-4 most common
		\item Each radix level trades fewer multiplications for code complexity and cache performance
		\item Modern FFT libraries are extremely complicated and performance is very hardware dependent 
		\item Not discussed here: in place? bit reversal for high-radix?
		\item Not discussed here: decimation in time / decimation in frequency
	\end{itemize}
\end{frame}

\begin{frame}[fragile]{Coarse Approach 5}
	
	\begin{lstlisting}[basicstyle=\ttfamily\scriptsize]
// N*log2(N) = 256*8 = 2048
void butterfly_dit(Complex* buffer) {
	// en.wikipedia.org/wiki/Cooley%E2%80%93Tukey_FFT_algorithm
	size_t m = 1;
	size_t exp_stride = Nsamples;
	for (size_t s = 1; s <= n_bits; ++s) {
		size_t mHalf = m;
		m *= 2;
		for (size_t k = 0; k < Nsamples; k+= m) {
			for (size_t j = 0; j < mHalf; ++j) {
				auto omega = twiddles[exp_stride*j];
				size_t indA = k + j;
				size_t indB = indA + mHalf;
				auto u = buffer[indA];
				auto t = omega * buffer[indB];
				buffer[indA] = u + t;
				buffer[indB] = u - t;
			}
		}
		exp_stride /= 2;
	}
}
	\end{lstlisting}
	
\end{frame}


\begin{frame}{Results}
	\coarsetable{1}{5}
\end{frame}


\begin{frame}
	Oh wait, computing twiddles with \texttt{std::exp} is less efficient and less accurate than \texttt{std::polar} since \texttt{std::exp} is generalized for complex powers and contains branches. cppreference says \texttt{std::polar} is ~4.5 times faster: https://en.cppreference.com/cpp/numeric/complex/polar
\end{frame}

\begin{frame}
\coarsetable{2}{5}
\end{frame}

\begin{frame}
	% performance on Teensy	
	\coarsetable{5}{5}
\end{frame}


